\icode{E-MG-SI-L3-261020-001}
如图所示，人造地球卫星 1 在圆形轨道\Rmnum{1}上运行，人造地球卫星 2 在椭圆轨道\Rmnum{2}上运行，其中椭圆轨道上的 \(A\) 点为远地点，\(B\) 点为近地点，两轨道相切于 \(A\) 点，下列说法正确的是\paren[]
	
\begin{tikzpicture}[font=\scriptsize]
	\draw[cyan,dashed] (.6,0) ellipse(1.4 and .8);
	\draw[dashed,teal] (0,0) circle(2);
	\fill (-.8,0)node[left]{\(B\)} circle(.04) (2,0)node[right]{\(A\)} circle(.04);
	\node at(135:2)[above left]{轨道\Rmnum{1}};
	\node at(.4,1)[above]{轨道\Rmnum{2}};
	\draw[thin,fill=blue!30] (-.4,0) circle(.26);
	\draw[thin] (-.14,0) arc(0:-180:.26 and .1);
\end{tikzpicture}
\begin{choices}
	\item 卫星 1 在轨道\Rmnum{1}上的速度大于\qty{7.9}{km/s}
	\item 卫星 1 在 \(A\) 点的加速度小于卫星 2 在 \(B\) 点的加速度
	\item 卫星 1 和卫星 2 在相同时间内与地球连线扫过的面积相等
	\item 卫星 1 在轨道\Rmnum{1}上的机械能大于在卫星 2 在轨道\Rmnum{2}上的机械能
\end{choices}
\begin{answer}
    
\end{answer}